% Generated draft: requires mathematical and source review.
\documentclass[11pt]{article}
\usepackage[margin=25mm]{geometry}
\usepackage{fontspec}
\setmainfont{Latin Modern Roman}
\usepackage{amsmath,amssymb,mathtools,unicode-math}
\setmathfont{Latin Modern Math}
\usepackage{enumitem,needspace,xurl,hyperref}
\setlength{\parindent}{0pt}
\setlength{\parskip}{5pt}
\setlength{\emergencystretch}{3em}
\newcommand{\N}{\mathbb{N}}
\newcommand{\Z}{\mathbb{Z}}
\newcommand{\Q}{\mathbb{Q}}
\newcommand{\R}{\mathbb{R}}
\newcommand{\C}{\mathbb{C}}
\newcommand{\F}{\mathbb{F}}
\newcommand{\A}{\mathbb{A}}
\newcommand{\PP}{\mathbb{P}}
\newcommand{\E}{\mathbb{E}}
\newcommand{\eps}{\varepsilon}
\newcommand{\dd}{\,\mathrm{d}}
\begin{document}
\section*{UnsolvedMath ID: 1100102}
\textbf{Source problem code:} catalogue entry
\subsection{Definitions and mathematical statement}
A finite $K(G,1)$ is a finite connected CW complex with fundamental group $G$ and contractible universal cover. A group admitting such a space is said to be of type $F$. A finitely generated group is Gromov hyperbolic if a Cayley graph for a finite generating set has $\delta$-thin geodesic triangles for some $\delta\geq0$.

Let $G$ admit a finite $K(G,1)$. Assume that $G$ has no subgroup isomorphic to $\Z\times\Z$, and that it is balanced in the following precise sense:
\[
\forall x\in G\ \bigl(\operatorname{ord}(x)=\infty\bigr),\quad
\forall m,n\in\Z,\quad
\bigl(\exists g\in G:\ gx^mg^{-1}=x^n\bigr)
\Longrightarrow |m|=|n|.
\]
Bestvina's Question 1.2 asks whether every such $G$ is Gromov hyperbolic. There is no assumption that $G$ is already a subgroup of a hyperbolic group; that is a property of the reported counterexample.

The supplied research summary reports a negative answer by Italiano--Martelli--Migliorini: a nonhyperbolic group $G$ of type $F$, with geometric and cohomological dimension $4$, that embeds in a hyperbolic group. It reports that this embedding implies the absence of $\Z^2$ and the displayed balanced property, the latter by translation lengths. It also reports a simultaneous negative answer to Bestvina's Question 1.1 concerning the absence of Baumslag--Solitar subgroups. The described construction uses circle-valued Morse theory and Bestvina--Brady-type finiteness arguments on fibering cusped hyperbolic $5$-manifolds, with a fiber-kernel after suitable filling. These are statements of the supplied literature assessment, whose exact construction and hypotheses have not been independently checked in this draft.
\subsection{Short English statement}
Bestvina asks whether a group with a finite classifying space must be hyperbolic if it contains no $\Z^2$ and conjugate powers of any infinite-order element have equal absolute exponents. The supplied research summary reports a nonhyperbolic type-$F$ subgroup of a hyperbolic group as a counterexample.
\subsection{Sources}
\begin{itemize}
\item Local UnsolvedMath registry, permanent ID 1100102, exact source code catalogue entry, ``Questions in Geometric Group Theory --- Q 1.2''; registry status solved, research classification SOLVED-IN-LITERATURE.
\item Bestvina, \emph{Questions in Geometric Group Theory} (2004), Question 1.2, PDF page 2, \url{https://www.math.utah.edu/~bestvina/eprints/questions-updated.pdf}. The record identifies a manual transcription of a public author-hosted PDF, accessed July 29, 2026. It states that the rendered source presents the question as open and that current status was not independently verified.
\item Italiano--Martelli--Migliorini, \emph{Inventiones Mathematicae} 231 (2023), 1--38; arXiv:2105.14795. The supplied research summary attributes the counterexample and construction described above to this work. A paper title and precise theorem number are not supplied.
\end{itemize}
\end{document}
